\chapter{卷积、数据流与算子融合}

卷积是视觉模型中的核心算子。它的数学形式和矩阵乘不同，但在实现上常常与 GEMM 有密切关系。理解卷积性能，不能只看卷积核大小，还要看输入布局、输出布局、stride、padding、通道数、batch、缓存复用和是否需要中间展开。

\section{卷积的循环结构}

以二维卷积为例，每个输出位置会遍历输入通道和卷积核窗口。若输入形状为 $N \times C \times H \times W$，输出通道为 $K$，卷积核为 $R \times S$，那么每个输出元素大约需要 $C R S$ 次乘加。输出位置很多，权重会在不同图像位置复用，输入像素也会被相邻窗口复用。

教学版直接卷积可以写成多重循环。代码看起来长，但每一层循环都有明确含义：batch、输出通道、输出高宽、输入通道、卷积核高宽。

\begin{lstlisting}[language=C++,caption={教学版 NCHW 直接卷积}]
#include <algorithm>
#include <cstdint>
#include <span>

void conv2d_nchw_direct(std::span<const float> input,
                        std::span<const float> weight,
                        std::span<float> output,
                        std::int32_t batch,
                        std::int32_t in_channels,
                        std::int32_t in_h,
                        std::int32_t in_w,
                        std::int32_t out_channels,
                        std::int32_t kernel_h,
                        std::int32_t kernel_w,
                        std::int32_t out_h,
                        std::int32_t out_w) {
    std::fill(output.begin(), output.end(), 0.0F);
    for (std::int32_t n = 0; n < batch; ++n) {
        for (std::int32_t oc = 0; oc < out_channels; ++oc) {
            for (std::int32_t oh = 0; oh < out_h; ++oh) {
                for (std::int32_t ow = 0; ow < out_w; ++ow) {
                    float acc = 0.0F;
                    for (std::int32_t ic = 0; ic < in_channels; ++ic) {
                        for (std::int32_t kh = 0; kh < kernel_h; ++kh) {
                            for (std::int32_t kw = 0; kw < kernel_w; ++kw) {
                                const std::int32_t ih = oh + kh;
                                const std::int32_t iw = ow + kw;
                                const std::size_t input_index =
                                    ((static_cast<std::size_t>(n) * in_channels + ic) *
                                     in_h + ih) * in_w + iw;
                                const std::size_t weight_index =
                                    ((static_cast<std::size_t>(oc) * in_channels + ic) *
                                     kernel_h + kh) * kernel_w + kw;
                                acc += input[input_index] * weight[weight_index];
                            }
                        }
                    }
                    const std::size_t output_index =
                        ((static_cast<std::size_t>(n) * out_channels + oc) *
                         out_h + oh) * out_w + ow;
                    output[output_index] = acc;
                }
            }
        }
    }
}
\end{lstlisting}

这段代码没有处理 padding 和 stride，是为了突出主干。真实卷积会在边界判断、布局选择和向量化上更复杂。若每个输出位置都重新读取卷积窗口，输入复用没有充分发挥；若输出通道较多，同一个输入窗口可以同时用于多个输出通道，类似矩阵乘中 $A$ 元素复用于多个 $C$ 元素。

\section{im2col：用内存换 GEMM}

一种经典方法是 \term{im2col}{image to column}。它把每个卷积窗口展开成一列或一行，把卷积转化为矩阵乘。好处是可以复用高度优化的 GEMM；坏处是展开后的中间矩阵可能很大，带来额外内存流量和空间占用。im2col 的思想非常重要，因为它展示了一个常见工程取舍：把复杂访问变成规则 GEMM，但付出数据重排成本。

当 batch 大、通道多、输出位置多时，GEMM 的高效率可能覆盖 im2col 成本；当模型小、内存紧、低延迟要求高时，直接卷积或隐式 GEMM 可能更合适。工业库经常根据形状选择不同算法，而不是固定一种实现。

\section{Winograd 和特殊卷积}

对于小卷积核，尤其是 $3 \times 3$，Winograd 算法可以减少乘法数量，但会增加加法、变换和数值误差风险。深度可分离卷积则把空间卷积和通道混合拆开，显著减少操作数，但也降低计算强度，使内存访问和布局更重要。$1 \times 1$ 卷积本质上接近矩阵乘，常常直接走 GEMM 路径。

这些变种说明卷积没有单一最优路径。算法层面减少乘法，不一定总能加速；若减少计算后程序变成内存受限，收益会被带宽限制吃掉。选择卷积算法时要同时看形状、数据类型、缓存、SIMD 和误差容忍。

\section{算子融合}

卷积后面常跟 bias、BatchNorm、ReLU、残差加法等操作。若每个操作都把完整中间张量写回内存，再由下一个操作读出来，内存流量会很大。算子融合把多个操作合并到同一个内核里，例如在卷积累加完成后立即加 bias 并做 ReLU，只把最终结果写回。

\begin{lstlisting}[language=C++,caption={卷积输出阶段融合 bias 和 ReLU 的形状}]
float finish_conv_value(float acc, float bias) {
    const float with_bias = acc + bias;
    return with_bias > 0.0F ? with_bias : 0.0F;
}
\end{lstlisting}

融合的边界不是越多越好。融合过多会让内核复杂、寄存器压力上升、代码生成路径爆炸，也会减少中间结果复用的灵活性。合理融合通常针对读写流量大的简单逐元素操作，让它们跟生产中间值的重算子靠在一起。

\begin{keyidea}
卷积优化的主线是数据流选择：直接卷积少中间数据但访问复杂，im2col 访问规则但增加展开流量，Winograd 减乘法但增加变换和误差，融合减少中间张量但增加内核复杂度。
\end{keyidea}

\begin{exercisebox}
选择一个 $3 \times 3$ 卷积形状，估算直接卷积的 FLOPs 和输入、权重、输出字节数。再估算 im2col 展开矩阵的大小。比较二者的额外内存成本，并解释为什么库需要按形状选择算法。
\end{exercisebox}
